\chapter{Introduction}
\label{ch:introduction}

\section{Background and Motivation}

Serverless computing, particularly Function-as-a-Service (FaaS), represents a significant evolution in cloud computing paradigms. By abstracting infrastructure management entirely, serverless platforms enable developers to deploy event-driven, auto-scaled functions without provisioning or managing servers \citep{shafiei2022survey}. AWS Lambda, the dominant FaaS platform, charges only for actual compute time (billed per 100 ms) plus a per-request fee, eliminating costs during idle periods.

However, this economic model introduces a critical performance trade-off: \textit{cold-start latency}. When a function has not been invoked recently, the platform must initialize a new execution environment, a process involving micro-VM creation, runtime bootstrap (Python interpreter, Node.js V8 engine, or Java Virtual Machine), and application code loading including all dependencies \citep{golec2024coldstart}. This initialization overhead, termed \textit{Init Duration} in AWS CloudWatch terminology, can range from 100 milliseconds to several seconds, severely degrading user-perceived performance for synchronous APIs and cascading microservice architectures.

\citet{eismann2022state} report that 48\% of serverless practitioners identify cold starts as a primary concern, with tail latencies (P95, P99) disproportionately affected. For latency-sensitive applications such as real-time payment processing, recommendation engines, or IoT event handlers, cold-start delays can violate service-level objectives (SLOs) and degrade user experience.

\section{Problem Statement}

While platform-level innovations such as VM snapshotting \citep{ao2022faasnap} and checkpoint-restore mechanisms \citep{kohli2024pronghorn} demonstrate significant cold-start reduction, these techniques are either unavailable on AWS Lambda or restricted to specific runtimes (e.g., SnapStart for Java 11/17 only). Provisioned concurrency, AWS's official mitigation strategy, eliminates cold starts by maintaining pre-initialized containers continuously, but incurs baseline costs proportional to provisioned capacity, contradicting the pay-per-use serverless economic model \citep{shafiei2022survey}.

Developers require cost-effective, runtime-agnostic optimization strategies using only standard AWS Lambda APIs. While prior work investigates individual controls—runtime selection \citep{wen2023characterizing}, memory allocation \citep{bluemke2025ijet}, or dependency management \citep{dantas2022deployment}—no systematic study compares these controls with explicit \textit{Init Duration} isolation across multiple production runtimes.

\citet{bluemke2025ijet} establish a rigorous methodology for AWS Lambda performance-cost analysis, evaluating memory configurations and processor architectures (x86\_64 vs. ARM64 Graviton2). However, their measurements aggregate cold-start initialization and handler execution into a single Duration metric, preventing identification of where latency originates. Additionally, their study examines only Python 3.9 and does not investigate deployment package characteristics (dependency bloat vs. minimal configurations).

This research addresses these limitations by isolating Init Duration and systematically comparing developer-adjustable controls across three production runtimes.

\section{Research Question}

The central inquiry guiding this investigation is:

\begin{quote}
\textit{How much of AWS Lambda cold-start latency can be removed by developer-adjustable controls, and at what cost?}
\end{quote}

Developer-adjustable controls in scope are:

\begin{enumerate}
    \item \textbf{Runtime language selection:} Python 3.12, Node.js 20, Java 21 (production-ready runtimes as of September 2026)
    \item \textbf{Deployment package optimization:} Bloated (unused dependencies included) vs. minimal (stdlib/no-dependency only)
    \item \textbf{Memory allocation:} 128 MB to 3008 MB (affects vCPU allocation above 1769 MB threshold)
    \item \textbf{Low-frequency warming:} EventBridge scheduled invocations (every 5 minutes) as a zero-baseline-cost alternative to provisioned concurrency
\end{enumerate}

Explicitly excluded: provisioned concurrency (paid continuous warm-up), SnapStart (Java-only, requires Corretto 11/17), custom AMIs, or platform modifications.

\section{Research Objectives}

This study pursues four specific objectives:

\begin{enumerate}
    \item \textbf{Isolate and quantify Init Duration} across three production runtimes (Python 3.12, Node.js 20, Java 21) to determine baseline cold-start magnitude per language.
    
    \item \textbf{Measure package-size effects} by comparing bloated deployments (including popular but unused dependencies) against minimal configurations (standard library or zero dependencies).
    
    \item \textbf{Evaluate low-frequency warming impact} on cold-start frequency by deploying EventBridge scheduled invocations (5-minute intervals) and measuring the proportion of subsequent invocations that encounter cold starts.
    
    \item \textbf{Express findings as a decision matrix} showing milliseconds saved per additional dollar per 1,000 invocations, enabling practitioners to select optimizations matching their latency and budget constraints.
\end{enumerate}

\section{Research Hypotheses}

Based on the literature review and preliminary proxy benchmarks (local process initialization measurements), we formulate four testable hypotheses:

\begin{itemize}
    \item \textbf{H1 (Runtime):} Init Duration differs significantly between Python, Node.js, and Java runtimes under minimal package configurations. Test: Kruskal-Wallis (if non-normal) or one-way ANOVA (if normal) with Holm-Bonferroni correction.
    
    \item \textbf{H2 (Package Size):} Deployment package pruning (removing unused dependencies) significantly reduces Init Duration within each runtime. Test: Mann-Whitney U (if non-normal) or Welch t-test (if normal) per runtime, with Holm-Bonferroni correction over three comparisons.
    
    \item \textbf{H3 (Warming):} Low-frequency EventBridge warming significantly reduces cold-start frequency compared to unwarm baseline. Test: Paired t-test or Wilcoxon signed-rank on 20-minute blocks, plus Fisher's exact test on cold-start proportions.
    
    \item \textbf{H4 (Memory, Exploratory):} Memory allocation impacts Init Duration due to vCPU coupling above 1769 MB. Test: Kruskal-Wallis or one-way ANOVA with post-hoc pairwise comparisons.
\end{itemize}

All hypothesis tests are pre-registered in the analysis plan (\texttt{docs/ANALYSIS\_PLAN.md}) with explicit test selection criteria (Shapiro-Wilk normality at $\alpha = 0.05$), multiple comparison correction (Holm-Bonferroni), and mandatory effect size reporting (epsilon-squared, rank-biserial, Cohen's d, or eta-squared as appropriate).

\section{Significance and Contributions}

This research makes the following contributions:

\begin{enumerate}
    \item \textbf{Methodological:} Extends \citet{bluemke2025ijet}'s configuration analysis framework with explicit Init Duration isolation, enabling precise identification of initialization overhead vs. execution time.
    
    \item \textbf{Empirical:} First systematic comparison of developer-adjustable cold-start controls across three production runtimes with controlled package-size conditions.
    
    \item \textbf{Practical:} A decision matrix expressing optimizations as \textit{cost-benefit ratios} (ms saved per dollar per 1,000 invocations), directly actionable by practitioners.
    
    \item \textbf{Reproducible:} Fully-instrumented Infrastructure-as-Code (AWS SAM template), standardized workload (deterministic SHA-256 computation), comprehensive test suite (100+ pytest tests), and data provenance tracking (live/mock/proxy mode segregation).
\end{enumerate}

\section{Document Structure}

The remainder of this report is organized as follows:

\begin{itemize}
    \item \textbf{Chapter 2 (Literature Survey):} Reviews state-of-the-art in serverless cold-start characterization, optimization strategies, and benchmarking methodologies, identifying the research gap this study addresses.
    
    \item \textbf{Chapter 3 (Research Methodology):} Details the experimental design, including controlled variables, measurement instrumentation, statistical analysis plan, and ethics considerations.
    
    \item \textbf{Chapter 4 (Design Specifications):} Describes the architecture of the Lambda functions, Infrastructure-as-Code implementation, workload design, and data collection pipeline.
    
    \item \textbf{Chapter 5 (Implementation):} Presents packaging strategies, deployment procedures, log parsing mechanisms, and cost modeling formulas.
    
    \item \textbf{Chapter 6 (Evaluation):} Reports experimental results, statistical analysis outcomes, and decision matrix construction (contingent on live AWS execution).
    
    \item \textbf{Chapter 7 (Conclusions):} Summarizes findings, discusses implications for practitioners, acknowledges limitations, and outlines future research directions.
\end{itemize}

This chapter established the research context, problem statement, objectives, and hypotheses. The next chapter reviews relevant literature to position this work within the broader serverless computing research landscape.